\documentclass[11pt]{article}

\usepackage[T1]{fontenc}
\usepackage[utf8]{inputenc}
\usepackage{lmodern}
\usepackage[margin=1in]{geometry}
\usepackage{microtype}
\usepackage[numbers]{natbib}
\usepackage{hyperref}
\usepackage{enumitem}

\hypersetup{
  colorlinks=true,
  linkcolor=black,
  citecolor=blue,
  urlcolor=blue
}

\setlength{\parskip}{0.7em}
\setlength{\parindent}{0pt}
\setlist[itemize]{leftmargin=1.4em}
\setlist[enumerate]{leftmargin=1.6em}
\emergencystretch=3em

\title{\textbf{OpenScientists:\\Science Is the Source Code of Civilization.\\It Must Not Be Compiled Behind Closed Doors.}}
\author{}
\date{April 2026}

\begin{document}
\maketitle

\begin{abstract}
Every technology that ever mattered---fire, writing, the printing press, the transistor---mattered because it changed what humans could \emph{know} and how fast they could act on that knowledge. AI-assisted scientific research is not another entry on that list. It is a discontinuity. For the first time, the engine that produces all downstream technologies---science itself---is becoming partially automatable. Hypothesis generation, experimental design, evidence synthesis, cross-domain reasoning: the cognitive core of discovery is now subject to the same compounding acceleration that Moore's Law once applied to circuits. If this transition is executed well, it compresses the timeline on every hard problem civilization faces---energy, disease, materials, climate. If it is executed poorly, or captured by a handful of private platforms, the most consequential knowledge-production system in human history will be shaped by incentives that have nothing to do with truth. This is not a policy question or a market question. It is \emph{the} question: who controls the rate at which humanity learns?

This paper proposes \emph{OpenScientists}, a public-interest initiative to seize that question before it is answered by default. We describe five projects---a structured library of researcher tacit knowledge, an agent runtime for scientific workflows, a benchmark grounded in real research practice, model-development work tuned for scientific reasoning, and a metascience effort studying how AI reshapes scientific production. Together they form a single research program aimed at ensuring that the infrastructure of AI-assisted discovery is open, inspectable, and governed by scientific norms rather than commercial strategy. The window is closing. The defaults are being set now.
\end{abstract}

\section{Why This Is the Most Important Transition in Human History}

\subsection{Science is the meta-technology}

There is a hierarchy among technologies. Most technologies solve particular problems: the internal combustion engine moves vehicles, antibiotics kill bacteria, semiconductors compute. But science is not a technology in this sense. Science is the \emph{process} that generates technologies. It is the meta-technology, the upstream engine whose output is every other capability civilization possesses. Agriculture, medicine, engineering, computation---all of them are downstream of the scientific method applied to the physical world.

This distinction matters enormously for understanding what it means to automate science. When you automate manufacturing, you produce more goods. When you automate logistics, you move goods faster. When you automate science, you accelerate the process that produces \emph{every future technology}. There is no leverage point in human civilization that is higher.

And yet the machinery is straining. Bloom et al.\ show that sustaining a given rate of productivity growth now requires far more research effort than it did a few decades ago \citep{bloom2020ideas}. The problems are harder. The knowledge frontiers are further apart. The coordination costs are immense. Discovery is not scaling with the problems it needs to solve.

\subsection{AI changes the equation}

AI has already compressed individual research steps. AlphaFold did not merely advance protein structure prediction---it effectively \emph{closed} a fifty-year grand challenge and made its output a routine tool for thousands of labs worldwide \citep{jumper2021alphafold}. Broader surveys describe advances in representation learning, hypothesis generation, simulation, and experiment design across fields \citep{wang2023discovery}. These are not parlor tricks. They are early demonstrations that large parts of the cognitive work of research---the scanning, the synthesis, the pattern recognition, the hypothesis formation---are computable.

No one has built a general AI scientist. But the trajectory is unmistakable. The question is no longer \emph{whether} AI will reshape scientific research. The question is who builds the infrastructure, who sets the defaults, and whose interests the system serves.

\subsection{What makes discovery hard---and why AI hits the right bottleneck}

The bottleneck in modern science is not intelligence. It is coordination. Hypothesis spaces are vast. Relevant evidence is scattered across papers, datasets, codebases, instruments, and---crucially---the tacit judgment of experienced researchers that never appears in any publication. Experiments are expensive. Negative results go unreported. Cross-disciplinary translation is agonizingly slow. The raw ingredients of a breakthrough often exist, scattered across labs and literatures, for years or decades before anyone assembles them.

This is the exact shape of problem that AI systems are best positioned to attack: absorbing scattered evidence at scale, preserving local know-how, surfacing hidden connections, and proposing next steps under uncertainty. A strong researcher with AI support gains the throughput of a well-supported team. A small lab gains synthesis capacity that previously required elite institutional resources. And if these systems also capture and transmit the tacit skills of experienced researchers, they don't just accelerate science---they improve its \emph{institutional memory}, preserving know-how that currently lives in individual heads and vanishes with career transitions.

\section{The Open Infrastructure Imperative}

\subsection{The race is already running}

This is not a thought experiment. The race to build AI scientist infrastructure is underway, and it is being led by the same concentrated group of frontier AI labs that already dominate foundation models.

OpenAI launched \emph{deep research} in February 2025, framing multi-step knowledge synthesis as a step toward systems capable of novel scientific work \citep{openai2025deepresearch}. Google described \emph{AI co-scientist} in March 2025, with a technical report on multi-agent biomedical hypothesis generation \citep{google2025health,gottweis2025coscientist}. Anthropic detailed its production \emph{multi-agent research system} in June 2025, emphasizing parallel exploration and structured citation support \citep{anthropic2025researchsystem}.

The direction is unambiguous. The investment is massive. And every month that passes without a credible open alternative is a month in which closed systems accumulate users, workflows, institutional habits, and lock-in.

\subsection{This is an epistemic crisis, not just a market problem}

The standard case for open source emphasizes access and cost. Those matter. But they radically understate what is at stake.

When a small number of corporations build the systems that mediate hypothesis generation, evidence synthesis, and experimental design, the issue is not merely who can afford the subscription. It is who shapes the \emph{epistemic defaults} of science itself. Which literatures get retrieved. Which reasoning patterns get rewarded. Which failure modes stay hidden inside proprietary pipelines. Which questions get asked and which get quietly deprioritized because they don't serve commercial objectives.

The compute divide between industry and academia has already widened to a chasm \citep{besiroglu2024computedivide}. Scale economies in foundation models produce durable concentration \citep{korinek2024concentrating}. If scientific discovery itself becomes mediated by frontier AI, then concentration in AI becomes concentration over the \emph{direction of human knowledge}. This is not an economic externality. It is the most dangerous epistemic bottleneck in the history of organized research.

A public counterweight is not optional. It is existential.

\begin{itemize}
  \item \textbf{Legitimacy.} Science built on opaque systems cannot be trusted, reproduced, or contested. Open infrastructure preserves the norms that make scientific claims worth believing.
  \item \textbf{Reproducibility.} Shared benchmarks, open protocols, and inspectable workflows are the only way the community can verify whether a claimed capability is real or a product of selective demonstration.
  \item \textbf{Pluralism.} Open infrastructure lets researchers adapt tools to their own questions, rather than inheriting a single commercial agenda as their analytical lens.
  \item \textbf{Contestability.} When a system produces a result---positive or negative---the community must be able to inspect the reasoning chain end-to-end. ``Trust us'' is not a scientific standard.
\end{itemize}

Nosek et al.\ argue that scientific credibility depends on the community's ability to review, critique, extend, and reproduce claims \citep{nosek2015openresearch}. Munaf\`{o} et al.\ emphasize that transparency, reproducibility, and aligned incentives are necessary for efficient knowledge accumulation \citep{munafo2017manifesto}. These principles are not weakened when AI enters the research process. They become \emph{non-negotiable}.

Open infrastructure is not naive openness. Not every capability must be released on the same timeline. Biosecurity, misuse, and low-quality automation are real concerns. But those concerns \emph{strengthen} the case for public-interest infrastructure and independent evaluation. The alternative is not safety. The alternative is unaccountable power over the direction of science, held by organizations whose primary obligation is to shareholders.

\section{The OpenScientists Mission}

\emph{OpenScientists} is a public-interest initiative to build the open infrastructure for AI-assisted scientific discovery before the window closes. \emph{OpenScientist} is one project within it, focused in part on capturing and operationalizing researcher know-how.

The mission is direct: build the full stack required for trustworthy, useful, and widely accessible AI scientist systems---not as a single model release, but as interlocking infrastructure:

\begin{itemize}
  \item A structured layer for capturing and reusing tacit researcher knowledge---the expertise that never appears in papers.
  \item A runtime for scientific agents with tool use, memory, and workflow support.
  \item Evaluations that measure genuine scientific ability, not proxy tasks.
  \item Model-development efforts tuned specifically for scientific reasoning.
  \item Governance and metascience work studying how these systems reshape research itself.
\end{itemize}

The strategic logic is simple and urgent: open infrastructure must exist \emph{before} closed systems become the default substrate for research. Once workflows, evaluations, and institutional habits lock in around proprietary providers, the cost of building a credible alternative rises by an order of magnitude. We are in the window now. It will not stay open.

Five operating principles:

\begin{enumerate}
  \item \textbf{Usefulness first.} Build systems that help with real research, not polished demos that impress on Twitter and fail in the lab.
  \item \textbf{Reproducibility by default.} Every method, dataset, and workflow must be inspectable and rerunnable by anyone.
  \item \textbf{Augmentation, not replacement.} The goal is to make every scientist dramatically more effective, not to render human expertise obsolete.
  \item \textbf{Respect field differences.} Biology, materials science, mathematics, and physics have radically different workflows. The infrastructure must serve all of them.
  \item \textbf{Public goods over lock-in.} Outputs must be portable, interoperable, and accessible to institutions far beyond elite labs.
\end{enumerate}

\section{Projects and Milestones}

These five projects form a single pipeline. The tacit-knowledge library supplies the domain grounding that agents need to behave like competent researchers rather than generic reasoners. The agent runtime provides the execution layer where that knowledge is deployed. The benchmark measures whether the combination actually works on real scientific tasks. Model development closes the loop by training on the signals the benchmark produces. And the metascience project monitors downstream effects on scientific production, feeding governance insights back into design decisions across the entire stack. Each layer depends on the ones before it and informs the ones after. This is not a collection of independent projects. It is one machine.

\subsection{A global library of researcher tacit knowledge}

Most of what makes a researcher effective never appears in a paper. It lives in heuristics, experimental instincts, debugging routines, failure cases, and the hard-won judgment about what usually breaks and why. This is the dark matter of scientific expertise. Without it, AI scientist systems are limited to the explicit record---competent at retrieval, incapable of the judgment calls that separate productive research from expensive busywork.

OpenScientists will build a structured, versioned, reviewable library of this tacit knowledge: triggers, decision rules, anti-patterns, failure modes, and domain-specific context drawn from working researchers across every field. This is not documentation. This is the construction of \emph{shared research memory}---know-how that would otherwise remain trapped in individual labs, disappear with career transitions, or die with retirements. Preserved, attributed, and made usable by both humans and agents, it becomes the most valuable dataset in the history of science.

This library is the foundation of the entire stack. Without it, everything downstream is a generic tool-use harness. With it, the system has judgment.

\subsection{An agent runtime and workflow system}

The second project is the execution layer: the infrastructure that lets AI scientist systems actually \emph{do work}. Literature search, structured planning, code execution, simulation, notebook interaction, persistent memory, citation tracing, review loops, and interfaces to domain-specific instruments and tools.

The runtime draws on the tacit-knowledge library to inform agent behavior. It must be modular enough for many fields but disciplined enough for rigorous evaluation. The goal is not one autonomous agent that does everything. The goal is a reliable, inspectable runtime where different models, prompts, tools, and knowledge configurations can be compared head-to-head on shared tasks. This testbed is what makes the benchmark possible.

\subsection{A benchmark for real scientific ability}

Most existing benchmarks in this space measure the wrong thing. Answering multiple-choice questions about papers is not science. Summarizing abstracts is not science. The tasks that actually matter---generating a credible hypothesis, spotting a confounder, designing an ablation study, identifying a hidden assumption, knowing when evidence is too weak to proceed---are almost entirely absent from current evaluations.

OpenScientists will develop a benchmark grounded in genuine research practice: difficult tasks designed with expert input, evaluated with expert review, and assessed on process quality, not just output fluency. It must reward caution, calibration, and error correction. Without serious evaluation, the entire field of AI-assisted science will be governed by anecdote, cherry-picked demos, and marketing decks. The benchmark is the immune system of the research program.

It also produces the signal that model development needs: which reasoning failures are systematic, which domains remain hardest, and where tacit-knowledge coverage has gaps.

\subsection{Model development for scientific reasoning}

General-purpose post-training will not be enough. Scientific reasoning demands specialized capabilities: structured tool use, long-horizon planning, methodological critique, calibrated uncertainty, and deep interaction with real research artifacts---notebooks, protocols, figures, code, reviews.

OpenScientists will support a staged model-development effort. Early phases focus on open datasets, synthetic tasks, distillation, retrieval, and post-training. Later phases include purpose-built foundation-model work if resources permit. Training signal comes from the benchmark. Domain grounding comes from the tacit-knowledge library. The point is not scale for its own sake but models whose training is \emph{aligned with scientific norms}---models that are honest about uncertainty, transparent about reasoning, and optimized for truth rather than fluent persuasion.

\subsection{Studying the effects on scientific production}

Any infrastructure this powerful must study its own failure modes. The most immediate risk is a flood of AI-generated manuscripts that raise reviewer burden, degrade signal quality, and crowd out genuine contributions. LLM usage in scientific writing is already rising at scale \citep{liang2024llmpapers}. Much of this is benign or productivity-enhancing. But without measurement, the line between assistance and pollution is invisible.

OpenScientists will run a dedicated metascience project: measuring prevalence, studying reviewer load, evaluating effects on citation networks, and testing interventions that preserve openness without letting scientific communication drown in noise. The findings feed directly back into governance decisions across the other four projects---what to release, how to release it, and what safeguards to build into the runtime. This is the feedback loop that keeps the entire initiative honest.

\section{Resources and Collaboration}

Building this infrastructure requires specific categories of people, compute, and institutional commitment. Nothing about this is optional. Every category below is load-bearing.

\subsection{Researchers contributing skills and tacit knowledge}

The knowledge library depends on PhD students and researchers across fields contributing the expertise that never makes it into publications: heuristics, failure cases, debugging routines, domain-specific judgment calls. These contributions are structured as versioned, reviewable skill artifacts---procedural rules, semantic corrections, and episodic case studies---organized by field and subfield. The value of the library scales directly with the breadth and depth of participation. A single lab's know-how is useful. The collective tacit knowledge of thousands of working researchers across disciplines is \emph{qualitatively different}---it is the difference between a notebook and a nervous system.

\subsection{Researchers as users of the system}

The same researchers who contribute skills must also be active users---running the AI scientist system in their own workflows, stress-testing it against real problems, and feeding back what works and what fails. This dual role is essential: contributors who also use the system can identify gaps in the knowledge base, catch skills that sound right but mislead in practice, and ground the infrastructure in actual research rather than synthetic benchmarks. Without this feedback loop, the system builds on assumptions rather than evidence.

\subsection{Professors and senior domain experts as reviewers and stewards}

Quality control requires experienced judgment. Professors and senior researchers serve as reviewers for contributed skills, stewards of domain-specific standards, and intellectual contributors who shape the research agenda. Their role is to ensure that the knowledge base meets the quality threshold of the fields it claims to represent---rejecting contributions that are vague, incorrect, or trivially obvious, and elevating those that encode genuinely hard-won insight. Without senior oversight, the knowledge library risks becoming a repository of confident mediocrity.

\subsection{Financial sponsors}

Sustained infrastructure development requires funding beyond what volunteer effort and academic grants can provide. Philanthropic sponsors, venture funds, and mission-aligned organizations can fund engineering teams, fellowships for contributing researchers, benchmark prizes, governance work, and the long-term maintenance that open-source projects require but almost never receive. The funding model must preserve independence. The initiative serves science, not its funders' portfolios.

\subsection{Token and API credit sponsors}

Evaluation, model iteration, and large-scale benchmarking require ongoing access to frontier language models. API credit sponsorships from model providers or intermediaries supply this access without requiring OpenScientists to build or host every model internally. This is particularly critical for the benchmarking and agent-runtime projects, where the entire point is comparing multiple models on shared tasks under controlled conditions.

\subsection{GPU sponsors}

Pre-training experiments, fine-tuning runs, large-scale evaluations, and agent simulations all require substantial compute. GPU sponsorships---from cloud providers, hardware companies, or national compute programs---are essential for the model-development and benchmarking tracks. Without dedicated compute, the open effort cannot credibly iterate on models or run evaluations at the scale needed to produce results that anyone takes seriously. Compute is not a nice-to-have. It is the currency of credibility in this field.

\subsection{Academic collaborators and operational contributors}

Universities, research institutes, scientific publishers, and professional societies can contribute in ways that go beyond individual participation: partnering on benchmark design, hosting pilot deployments, running policy experiments around AI-assisted research norms, contributing to reviewer-load studies, and providing institutional legitimacy. Operational contributors---project managers, community organizers, documentation writers---keep the initiative functional as it scales beyond a handful of founding participants.

\medskip

These seven categories are not a wishlist. They are the \emph{minimum viable coalition}. Researchers generate the knowledge. Reviewers ensure its quality. Sponsors provide the compute, funding, and API access that make iteration possible. Academic partners ground the work in institutional reality. Each mode of contribution reinforces the others. The initiative either assembles all of them or it does not succeed.

\section{Conclusion}

We are living through the most consequential transition in the history of organized knowledge. Frontier AI labs are building systems that will reshape scientific practice within years, not decades. The trajectory is clear: hypothesis generation, experimental design, evidence synthesis, and cross-domain reasoning are being automated, piece by piece, inside closed platforms operated by a handful of corporations.

The question is not whether this will happen. It is already happening. The question is whether the infrastructure mediating this transformation will be inspectable, contestable, and governed by scientific norms---or whether it will default to proprietary platforms whose incentives have nothing to do with the pursuit of truth.

OpenScientists is a proposal for the former: public-interest infrastructure that preserves scientific legitimacy while dramatically expanding scientific capacity. Five projects---tacit-knowledge library, agent runtime, benchmark, model development, metascience---form a single research program. Each depends on the others. Together they constitute the minimum viable open alternative to a future in which the direction of human knowledge is determined by commercial strategy.

Science is the source code of civilization. It must not be compiled behind closed doors. The time to build is now, because the defaults are being set as you read this sentence.

\bibliographystyle{plainnat}
\bibliography{refs}

\end{document}
